\section{Introduction}\label{sec:introduction}

A query on finite relational structures is a property unchanged by
isomorphism. It belongs to polynomial time if a deterministic algorithm
decides it in time polynomial in the size of an ordinary encoding of the
structure. The algorithm may use the order of that encoding; its answer
must be independent of the order. An unordered structure, in contrast,
does not supply such an order as one of its relations. Descriptive
complexity asks which formalisms can express every polynomial-time query
without first choosing an order.

Choiceless polynomial time with counting, denoted $\CPT$ throughout this
paper, operates on the hereditarily finite sets built from the input
elements as atoms. These are finite sets whose members are atoms or
further such sets, with no infinite descending membership chain. The
formalism permits set operations, comprehension and the cardinality
operation $\mathsf{Card}$, which returns a finite von Neumann ordinal.
Each fixed program has polynomial bounds on its computation and on the
number of recursively occurring objects. Its operations commute with
input automorphisms and do not provide an arbitrary choice of an ordering.
We use the established equivalent formulation by interpretation programs;
its exact syntax, counting operation and quotient semantics are stated in
Section~\ref{sec:coding}. In particular, $\CPT$ here always includes
counting and has no imposed bound on the rank of its sets.

To say that this formalism \emph{captures polynomial time} on unordered
finite structures means that every isomorphism-invariant polynomial-time
query, in every fixed finite relational vocabulary, is definable by a
program of the formalism. We first specify a query on all inputs in a
fixed vocabulary.

\subsection{The separating query}\label{sec:query}

Throughout, let $\mathbb F=\mathbb F_3$.  Fix the binary relational vocabulary
\[
 \tau=\{\mathsf{Ed},\mathsf{Cf},\mathsf{EB},\mathsf{VB},
             \mathsf I,\mathsf Z_0,\mathsf Z_1,\mathsf Z_2\}.
\]
The subscripts on the last three relation symbols denote elements of
$\mathbb F$.  For a finite $\tau$-structure $A$, set
\[
 Y=\{y\in A:\mathsf{Ed}(y,y)\},\qquad
 X=\{a\in A:\mathsf{Cf}(a,a)\},\qquad
 X_t=\{a\in X:\mathsf{VB}(t,a)\}\quad(t\in X).
\]
No disjointness or equivalence-relation assumption is imposed here.  In
particular, $X$ and $Y$ may overlap.  Define, in $\mathbb F$,
\begin{equation}\label{eq:query-coefficient}
 C(y,a)=\sum_{\substack{x\in Y\\\mathsf I(a,x)}}
          \ \sum_{\substack{\delta\in\mathbb F\\
                              \mathsf Z_\delta(y,x)}}\delta
       \qquad(y\in Y,\ a\in X).
\end{equation}
If several of the relations $\mathsf Z_\delta(y,x)$ hold, all their
contributions are included.

Introduce two separate families of unknowns,
$(\lambda_a)_{a\in X}$ and $(\mu_y)_{y\in Y}$.  Thus even when $a=y$,
$\lambda_a$ and $\mu_y$ are different unknowns.  The system defining our
query has the equations
\begin{equation}\label{eq:query-normalization}
 \sum_{a\in X_t}\lambda_a=1\qquad(t\in X)
\end{equation}
and, for every $(t,y)\in X\times Y$ for which
\begin{equation}\label{eq:query-incidence-test}
 \text{there exist $a\in X_t$ and $x\in Y$ with
          $\mathsf I(a,x)$ and $\mathsf{EB}(y,x)$,}
\end{equation}
the equation
\begin{equation}\label{eq:query-consistency}
 \mu_y=\sum_{a\in X_t}\lambda_a C(y,a).
\end{equation}
Let $Q(A)$ mean that this system is consistent over $\mathbb F$.
Repeated identical equations are harmless; they remain equations and are
not added together.  An empty sum is zero, and a system with no equations
is consistent.  These conventions make $Q$ a query on all finite
$\tau$-structures.

\subsection{Main result}

\begin{theorem}\label{thm:main}
The query $Q$ on all finite structures in the fixed vocabulary $\tau$ of
eight binary relations is isomorphism-invariant and decidable in
deterministic polynomial time, but is not definable in full choiceless
polynomial time with counting.
\end{theorem}

Theorem~\ref{thm:main} confirms the noncapture conjecture of Blass,
Gurevich and Shelah for the full counting formalism; its historical
formulation is discussed below. We first record the ordinary algorithmic
upper bound. The rest of the paper proves the choiceless lower bound.

\begin{proposition}\label{prop:query-ptime}
The query $Q$ is isomorphism-invariant and decidable in deterministic
polynomial time from an ordered encoding of its input structure.
\end{proposition}
\begin{proof}
An isomorphism carries $X,Y,X_t$ to their counterparts, preserves the
coefficient in \eqref{eq:query-coefficient}, and carries the test
\eqref{eq:query-incidence-test} to the corresponding test.  Consequently it
permutes the two families of unknowns and the equations compatibly.
Consistency is therefore invariant.

If the input has $n$ elements, there are at most $2n$ unknowns and
$n+n^2$ equations.  All coefficients and tests can be computed by
polynomially many loops over the input universe; for example, a direct
implementation uses $O(n^4)$ operations over the fixed field.  Gaussian
elimination then decides consistency in polynomial time
\cite[Sections~14.4--14.5]{shoup2009}. Any input order
may be used to list the rows and columns: changing it does not change
consistency.
\end{proof}

The same query is expressible in two standard linear-algebraic logics.
Following Gr\"adel and Pakusa \cite[Section~2]{graedelPakusa2019}, write
$\mathrm{FPS}_3$ for fixed-point logic with counting (FPC) extended by
the solvability operator over $\mathbb F_3$, and $\mathrm{FPR}_3$ for
FPC extended by the rank operator over $\mathbb F_3$. Their two-sorted syntax
allows bounded number variables alongside the unordered input elements.
We state the consequence for nonempty structures, so that the matrix
index sets in its proof are nonempty.

\begin{corollary}\label{cor:linear-algebraic-logics}
On nonempty finite $\tau$-structures, $Q$ is definable in both
$\mathrm{FPS}_3$ and $\mathrm{FPR}_3$. Consequently neither of these
logics is contained in $\CPT$.
\end{corollary}
FPC defines the coefficient predicates and the row test for $Q$.
Section~\ref{sec:linear-algebraic-logics} gives the resulting solvability
formula and compares the ranks of the coefficient and augmented matrices.
The nonempty witnesses constructed below give both noncontainments.

\subsection{Background and the obstruction to rank-dependent arguments}
The ordered case provides the starting point. Immerman and Vardi showed
that first-order logic with a least fixed-point operator expresses exactly
the polynomial-time queries on finite structures supplied with a linear
order~\cite{immerman1982,vardi1982}.
For unordered structures, Chandra and Harel asked whether polynomial-time
queries admit an effective enumeration~\cite[Section~5, p.~118]{chandraHarel1982}.
Gurevich formulated the question in terms of a logic capturing polynomial
time and conjectured that no such logic exists~\cite[Section~7]{gurevich1988}.
Excluding CPT with counting as a candidate leaves this broader question
open.

Blass, Gurevich and Shelah introduced choiceless polynomial time to study
computation on unordered structures using hereditarily finite sets rather
than arbitrary choices of an ordering~\cite{bgs1999}. Their original
1997 preprint explicitly conjectured that the extension by counting is
still a proper fragment of polynomial time~\cite[1997 version,
Introduction, p.~3]{bgs1999}. Their later study of unordered computation
examined CFI instances and finite-field linear algebra as possible sources
of separation, carefully distinguishing representations with common and
independent row and column index sets~\cite[Sections~4 and~6]{bgs2002}.
The query here is a particular linear-consistency problem over
$\mathbb F_3$, defined without a promise on the input structure.

Shelah's early work claimed noncapture for counting extensions of
choiceless computation~\cite[Introduction and Section~4]{shelah2000}.
Later accounts continued to state the capture problem as open, from
Blass, Gurevich and Shelah~\cite[Section~7]{bgs2002} through the 2025
work of Dawar, Gr{\"a}del, Kullmann and Pago~\cite[Section~2.2]{dawarGraedelKullmannPago2025}.
We give a separate proof for the explicit query above; Shelah's claim is
not a theorem input to this argument.

The construction of Cai, F{\"u}rer and Immerman established the central
example of a polynomial-time graph query beyond fixed-point logic with
counting~\cite{cfi1992}. Its local gadgets conceal a global parity
obstruction from bounded-variable counting formulas. This is an
antecedent of our charge-based witnesses, not an identification of those
witnesses with the original CFI graphs. The passage from counting logic
to choiceless computation is substantial: Dawar, Richerby and Rossman
showed that the preordered CFI query obtained from ordered base graphs is
definable without counting, while even counting does not suffice if set
rank is restricted to $o(\log n/\log\log n)$, where $n$ is the input
size~\cite[Theorems~17 and~40,
2007 author version]{drr2008}. In particular, their lower bound for every
fixed rank cannot by itself yield a separation for the full formalism.
Pakusa, Schalth{\"o}fer and Selman extended the positive result to
connected preordered base graphs whose colour classes have logarithmic
size in the number of base vertices~\cite[Theorem~1]{pakusaSchalthoeferSelman2016}.
The preorder orders classes of base vertices, rather than all vertices
of the resulting CFI structure.

Positive algebraic results also go beyond the original CFI family.
Abu Zaid, Gr{\"a}del, Grohe and Pakusa obtained choiceless canonization,
that is, construction of canonical ordered copies, when a supplied preorder has bounded classes whose induced substructures
have Abelian automorphism groups~\cite[Corollary~19]{abuZaidGraedelGrohePakusa2014}.
Their procedure uses specially structured cyclic systems of equations,
not arbitrary unordered linear systems~\cite[Definition~6 and
Theorem~14]{abuZaidGraedelGrohePakusa2014}. Lichter and Schweitzer obtained
analogous canonization for graphs when these induced groups are dihedral
or cyclic~\cite[Theorem~1]{lichterSchweitzer2021}. These results show why
the lower-bound argument must exploit the particular action of its
nonabelian group, rather than noncommutativity alone.

Functional lower bounds expose different limits. Rossman proved that
CPT with counting cannot construct the set of hyperplanes of an input
finite vector space over a fixed finite
field~\cite[Theorem~6.1, author version]{rossman2010}.
Pago proved that no such program constructs, on every hypercube with
$N$ vertices, a total preorder with classes of size
$O(\log N)$~\cite[Theorem~2 and Corollary~3, full version]{pago2021}.
The latter statement also excludes encodings from which the preorder can
be recovered in CPT. It obstructs the route of first constructing a fine
preorder and then applying the preordered CFI algorithm. These are
functional obstructions, not nondefinability of a Boolean query.

Pago's symmetric-XOR-circuit approach relates restricted classes of
choiceless CFI algorithms to circuit
families~\cite[Theorems~2, 31 and~33, full version]{pago2023}.
The translation permits polynomially many circuit images under
automorphisms of the base graph; the circuit lower bound assumes an
invariant circuit and different bounds on connections between gates.
These hypotheses differ, so that argument does not establish
nondefinability for the entire restricted algorithm class.

The support-and-transfer architecture has older roots. Blass, Gurevich
and Shelah developed supports and representations by forms evaluated at
supporting tuples, called molecules~\cite[Sections~8--9, published
version]{bgs1999}. Dawar, Richerby and Rossman developed rank-sensitive
support bounds and a counting transfer for hereditarily supported
objects~\cite[Sections~7--8, 2007 author version]{drr2008}. Using the
group-relative homogeneity from Section~\ref{sec:base-counting}, we adapt
that representation method in Section~\ref{sec:hf-transfer} and prove the
required exact counting statements.
The support estimate in Section~\ref{sec:support} instead uses a larger
nonabelian group to obtain small supports for its central subgroup,
independently of HF rank. Its standard class-two group and
alternating-form ingredients, as well as the projection and counting-game
methods used for the atom comparison, are credited at their respective
proofs. The grid-specific estimate and the quantitative transfer are
proved here.

Finally, the connection to arbitrary programs uses the interpretation
characterization of Gr{\"a}del, Pakusa, Schalth{\"o}fer and
Kaiser~\cite[Theorem~1]{gkps2015}. We use the binary, two-dimensional
formulation of Grohe, Schweitzer and Wiebking~\cite[Section~6, Theorem~18,
arXiv version~1]{gsw2020}, including its generated-equivalence and
existential-representative quotient convention. This characterization is
an external input. Section~\ref{sec:coding} proves the further facts
needed here: a polynomial bound on the complete hereditary encoding of a
run, and exact formulas for its states whose variable bound is independent
of the number of stages. Thus neither a bound on HF rank nor a restriction
to a particular algorithmic construction is imposed on the programs
excluded by the theorem.


\subsection{Proof overview}

For each side length $L$, we build two structures from the
three-dimensional box grid. Each edge carries states consisting of a
bounded collection of $\mathbb F_3^2$ face coordinates and one scalar.
A vertex configuration chooses incident edge states that agree on their
shared face coordinates and whose scalars have a prescribed signed sum,
called its charge. Both kinds of
states become input atoms, linked only by the fixed binary relations.
The two structures have all charges zero, or one charge equal to one.
Their size is $N=\Theta(L^3)$. The query's normalization equations assign
field-valued weights summing to one to the configurations at each vertex.
Its consistency equations make the weighted scalar value of an edge
agree at its two endpoints. Summing the vertex divergence equations then
forces the total charge to be zero, since each edge enters one vertex and
leaves another. This argument allows arbitrary field-valued weights;
it does not assume that a solution selects one configuration per vertex.
The all-zero configurations give a solution in the zero-charge structure,
so the query has opposite answers on the pair.

The challenge is to show that every fixed full $\CPT$ program eventually
gives the same answer on this pair. Three steps connect the symmetries of
these atoms to the complete computation.

\paragraph{Rank-independent supports.}
The structures have an automorphism group $H$ containing a central
subgroup $K$ of divergence-free scalar flows on the edges. A tuple of
atoms is a $K$-support of an object if every element of $K$ fixing that
tuple also fixes the object. For any fixed $q$, every hereditarily finite
object with $H$-orbit of size at most $N^q$ has a support of size
$O_q(L(1+\log L)^2)$. The bound does not depend on the object's rank.
Commutators in $H$ force the linear functionals detecting a stabilizer
to have nonzero curl on few faces. Subtracting explicitly constructed
gradients then concentrates those functionals on few edges. This is the
role of the nonabelian group; central flows alone do not supply this
argument.

\paragraph{Counting equivalence for all supported objects.}
For analysis only, we add a preorder identifying certain classes of
atoms and relations comparing scalar differences. These relations are
not given to the program. With
$M=\lfloor L^{7/4}\rfloor$, the two expanded structures agree on all
first-order counting sentences using at most $M$ variable names. A
charge discrepancy can be moved along paths in the large component that
remains after the edges touched by the assigned atoms are removed.
Within either expansion, equality of sufficiently wide counting types
also determines the $K$-orbit of a short tuple.

Fix a support length $s=O_q(L(1+\log L)^2)$. Consider all hereditarily
finite objects for which the object and every recursive member have a
$K$-support of length at most $s$. This is a countably infinite domain of
finite objects, containing objects of every finite rank. We represent
its objects by finite forms evaluated at supporting atom tuples, then
prove equality, membership and exact counting transfer. For each fixed
$m$, the two domains agree on all counting sentences with at most $m$
variable names when $L$ is sufficiently large. The key width cost is
$O_{q,m}(L^{3/2}(1+\log L)^3)=o(M)$.

\paragraph{The entire computation fits the comparison.}
The interpretation characterization supplies a fixed program whose
successive states are quotients of definable pairs from the preceding
state. We encode every quotient vertex by a tagged set of ordered pairs
of preceding vertex codes. A polynomial run has only polynomially many
such codes and recursive constituents; this transitive family is
$H$-invariant. The support theorem therefore places the entire encoded
trace in the domains just compared.

Finally, each state, halting test and output test has an exact
counting-formula definition over its supported domain. Its length may
grow with the input, but the number of variable names needed is bounded
by a constant depending only on the program. We prove this by a
capture-avoiding substitution lemma, including the generated quotient
relation and the cardinality comparisons. Choose the program and its
constants $q,m$ first, and then take $L$ large. Counting equivalence
forces the same halt stage and output on the two structures, contradicting
their opposite query answers.

\begin{figure}[t]
\centering
\begin{tikzpicture}[
  >=Stealth,
  box/.style={draw,rounded corners,align=center,font=\small,
              text width=3.65cm,minimum height=1.05cm,inner sep=5pt}
]
\node[box,text width=4.9cm] (input) at (0,0)
  {Two grid structures\\with zero or unit total charge};
\node[box] (query) at (-4.45,-1.65)
  {Linear-system query\\has opposite answers};
\node[box] (support) at (0,-1.65)
  {Small $H$-orbits give\\short $K$-supports};
\node[box] (base) at (4.45,-1.65)
  {Counting equivalence\\and homogeneity\\of expanded atoms};
\node[box] (trace) at (0,-3.4)
  {Trace codes in supported HF;\\bounded-variable\\state formulas};
\node[box] (hf) at (4.45,-3.4)
  {Counting equivalence\\on supported\\HF objects};
\node[box,text width=6.3cm] (equal) at (2.225,-5.3)
  {Each fixed CPT program\\has equal outputs for large $L$};
\draw[->] (input.south west)--(query.north);
\draw[->] (input.south)--(support.north);
\draw[->] (input.south east)--(base.north);
\draw[->] (support.south)--(trace.north);
\draw[->] (support.south east)--(hf.north west);
\draw[->] (base.south)--(hf.north);
\draw[->] (trace.south)--(equal.north west);
\draw[->] (hf.south)--(equal.north east);
\draw[<->,dashed] (query.south) |- 
  node[pos=.73,above,font=\small]{contradiction} (equal.west);
\end{tikzpicture}

\caption{The two routes from the input pair. The charge equations give
opposite query answers. The symmetry, logic and computation arguments
force every fixed choiceless program to give equal answers for sufficiently
large grids. The atom expansions are used only for analysis. If the
program decided $Q$, its equal outputs would contradict the left-hand route.}
\label{fig:proof-map}
\end{figure}

Section~\ref{sec:query} defines the query on all inputs, and
Section~\ref{sec:structures} constructs its two test structures.
Section~\ref{sec:support} proves the support theorem.
Sections~\ref{sec:base-counting} and~\ref{sec:hf-transfer} establish the
counting comparison on atoms and then on supported sets.
Section~\ref{sec:coding} encodes full interpretation computations, and
Section~\ref{sec:separation} fixes the parameters and proves
Theorem~\ref{thm:main}.
Section~\ref{sec:linear-algebraic-logics} proves the definability
of $Q$ in solvability and rank logic.
